% ============================================================================
%  packages.tex -- every package and every style decision for the template.
%
%  The numbers in here were measured off a reference thesis PDF, so don't
%  change them casually. Where a value looks odd there's a comment saying why.
% ============================================================================

% ---- encoding and fonts ----------------------------------------------------
%  times.sty only swaps the text families (rm -> Times, sf -> Helvetica,
%  tt -> Courier) and leaves maths in Computer Modern. That mix is deliberate:
%  it is what the reference document does. If you switch to mathptmx or newtx
%  the maths will change too and the look will drift.
\usepackage[utf8]{inputenc}
\usepackage[T1]{fontenc}
\usepackage{times}

% ---- language --------------------------------------------------------------
\usepackage[australian]{babel}
\usepackage{csquotes}

% ---- colour ----------------------------------------------------------------
\usepackage[dvipsnames,svgnames,table]{xcolor}

% ---- page geometry ---------------------------------------------------------
%  Text block is 16cm wide, 2.5cm from the top, 3cm off the bottom.
%  \headheight and \headsep stay at the report 11pt defaults (12pt / 25pt).
%  fancyhdr will warn that 12pt is a bit tight for the running head; ignore it.
%  Enlarging \headheight moves the header and breaks the match.
\usepackage[a4paper,left=2.5cm,right=2.5cm,top=2.5cm,bottom=3cm,headsep=26.61pt]{geometry}

% ---- line spacing ----------------------------------------------------------
%  1.1 x the 11pt default, i.e. 14.96pt baselines. Captions are pulled back to
%  single spacing further down, in the caption setup.
\usepackage{setspace}
\setstretch{1.1}

% ---- floats, graphics, tables ----------------------------------------------
\usepackage{graphicx}
\usepackage{booktabs}
\usepackage{array}
\usepackage{multirow}
\usepackage{longtable}
\usepackage{caption}
\usepackage{subcaption}
\usepackage{placeins}
\usepackage{rotating}
\usepackage{pdflscape}
\usepackage{tikz}
\usetikzlibrary{arrows.meta,positioning}

% ---- maths -----------------------------------------------------------------
\usepackage{amsmath}
\usepackage{amssymb}
\usepackage{mathtools}

% ---- lists -----------------------------------------------------------------
\usepackage{enumitem}

% ---- dummy text (delete this block once you have real writing) -------------
%  lipsum's Latin has no hyphenation patterns, so the odd paragraph would set
%  overfull without a little extra slack. \emergencystretch only kicks in on
%  lines that would otherwise overflow, so it changes nothing else.
\usepackage{lipsum}
\emergencystretch=1em

% ============================================================================
%  Headings
% ============================================================================
\usepackage{titlesec}

%  Chapter openers are two flush-left \huge bold lines, "Chapter N" then the
%  title, 8pt apart. (The report default puts the title at \Huge; this doesn't.)
\titleformat{\chapter}[display]
  {\normalfont\huge\bfseries}
  {\chaptertitlename\ \thechapter}{8.6pt}
  {\huge}
\titlespacing*{\chapter}{0pt}{50.6pt}{40.7pt}

%  Section sizes are exact point sizes rather than \Large / \large, because
%  that is what the reference measures at: 16pt, 14pt, 12pt.
\titleformat{\section}
  {\normalfont\fontsize{16}{18}\selectfont\bfseries}{\thesection}{1em}{}
\titleformat{\subsection}
  {\normalfont\fontsize{14}{16}\selectfont\bfseries}{\thesubsection}{1em}{}
\titleformat{\subsubsection}
  {\normalfont\fontsize{12}{14}\selectfont\bfseries}{\thesubsubsection}{1em}{}

\titlespacing*{\section}      {0pt}{3.5ex plus 1ex minus .2ex}{2.3ex plus .2ex}
\titlespacing*{\subsection}   {0pt}{3.25ex plus 1ex minus .2ex}{1.5ex plus .2ex}
\titlespacing*{\subsubsection}{0pt}{3.25ex plus 1ex minus .2ex}{1.5ex plus .2ex}

\setcounter{secnumdepth}{3}
\setcounter{tocdepth}{2}   % chapters + sections + subsections in the TOC

% ============================================================================
%  Running heads and page numbers
% ============================================================================
\usepackage{fancyhdr}
\pagestyle{fancy}
\fancyhf{}
%  Bold uppercase "N. CHAPTER TITLE", outer edge of the spread, thin rule
%  under it. Page number centred in the footer on every page.
\fancyhead[LE,RO]{\normalfont\bfseries\leftmark}
\fancyfoot[C]{\normalfont\thepage}
\renewcommand{\headrulewidth}{1pt}
\renewcommand{\footrulewidth}{0pt}
\renewcommand{\chaptermark}[1]{\markboth{\MakeUppercase{\thechapter.\ #1}}{}}
%  \chapter* (References, and anything else unnumbered) leaves the mark alone,
%  so set it by hand there -- see \unnumberedchapter in main.tex.

% ============================================================================
%  Captions
% ============================================================================
%  "Figure 2.3 Some caption text." -- bold label, no colon, a plain space
%  before the text. Single-line captions centre, longer ones justify to the
%  full text width. Captions are single spaced even though the body is 1.1.
\captionsetup{labelsep=space,labelfont=bf,font=singlespacing}

% ============================================================================
%  Cross-references and links
% ============================================================================
\usepackage{hyperref}
\usepackage{cleveref}
\usepackage{xurl}
\urlstyle{same}
\hypersetup{
  colorlinks=true,
  linktocpage=true,   % in the TOC/LoF/LoT the page number is the link, not the
                      % entry text, so the text stays black and only the number
                      % goes blue
  linkcolor=blue,     % TOC page numbers, \ref, page refs
  citecolor=teal,     % \cite
  urlcolor=blue,
  filecolor=blue,
}

% ============================================================================
%  Bibliography -- IEEE numeric, pdflatex + bibtex (NOT biber)
% ============================================================================
\usepackage[numbers,square,sort&compress]{natbib}
\bibliographystyle{IEEEtranN}
%  The "References" heading is set by \unnumberedchapter in main.tex, so stop
%  natbib adding a second "Bibliography" heading of its own.
\renewcommand{\bibsection}{}
